\section{Tổng quan nghiên cứu liên quan (Related Work)}
\label{sec:related_work}

\subsection{Các Mô hình Thống kê và Kinh tế lượng Truyền thống}
Dự báo phụ tải điện có lịch sử phát triển lâu dài với các phương pháp thống kê cổ điển như mô hình tự hồi quy tích hợp trung bình trượt (\textit{Autoregressive Integrated Moving Average} -- ARIMA), SARIMAX tích hợp biến ngoại sinh, và các mô hình san bằng mũ Holt-Winters~\cite{hyndman2018forecasting}. Các mô hình này sở hữu nền tảng toán học kinh tế lượng vững chắc, khả năng diễn giải nguyên nhân - kết quả tường minh và chi phí tính toán cực kỳ thấp. Tuy nhiên, hạn chế cố hữu của nhóm phương pháp này là giả định chặt chẽ về tính dừng (\textit{stationarity}) hoặc khả năng dừng sau khi lấy sai phân. Chúng hoàn toàn bất lực trong việc nắm bắt các mối tương tác phi tuyến tính bậc cao giữa chuỗi nhiệt độ, bức xạ mặt trời và sự biến động phức tạp của đồ thị phụ tải trong những ngày thời tiết dị thường.

\subsection{Học máy Dạng bảng và Gradient Boosted Decision Trees}
Sự trỗi dậy của các kỹ thuật học máy đã tạo nên cuộc cách mạng trong lĩnh vực STLF. Các công trình của Hong và Fan~\cite{hong2016probabilistic} khẳng định rằng các thuật toán cây quyết định tăng cường gradient (\textit{Gradient Boosted Decision Trees} -- GBDT), tiêu biểu là XGBoost~\cite{chen2016xgboost} và LightGBM~\cite{ke2017lightgbm}, luôn thể hiện ưu thế vượt trội trên các tập dữ liệu chuỗi thời gian dạng bảng có biến ngoại sinh. Ưu điểm nổi bật của GBDT bao gồm: khả năng nắm bắt tự nhiên các quan hệ phi tuyến và phi đơn điệu, khả năng chống chịu cao với các thang đo đặc trưng khác biệt mà không đòi hỏi chuẩn hóa khắt khe, và cơ chế cắt tỉa thông minh giúp kiểm soát phương sai.

Đặc biệt, công trình của Roy, Pyltsov và Hu~\cite{roy2025hourly} (arXiv:2505.11390) tại cuộc thi \textit{IISE PG\&E 2025} đã đề xuất chiến lược phân đoạn dữ liệu thành 24 bài toán con tương ứng với 24 giờ trong ngày và huấn luyện 24 mô hình XGBoost độc lập. Chiến lược này giúp triệt tiêu hoàn toàn sự biến thiên nhật triều mà không đòi hỏi các hàm ánh xạ phức tạp, đồng thời kết hợp PCA để giải quyết hiện tượng đa cộng tuyến giữa 5 trạm quan trắc nhiệt độ và 5 trạm bức xạ mặt trời. Kết quả thực nghiệm của bài báo gốc đã trở thành một cột mốc đối chuẩn quan trọng trong lĩnh vực dự báo phụ tải ngắn hạn.

\subsection{Deep Learning và Kiến trúc Transformers Chuỗi Thời gian}
Trong những năm gần đây, cộng đồng nghiên cứu đã nỗ lực áp dụng các mạng nơ-ron sâu (\textit{Deep Learning}), đặc biệt là kiến trúc \textit{Transformer} và cơ chế chú ý (\textit{Self-Attention}), vào bài toán dự báo chuỗi thời gian dài hạn và ngắn hạn. Nhiều kiến trúc tiêu biểu đã được giới thiệu như Informer~\cite{zhou2021informer} với cơ chế \textit{ProbSparse attention}, Autoformer~\cite{wu2021autoformer} kết hợp cơ chế phân tách chuỗi thời gian (\textit{series decomposition}) và khối tự tương quan (\textit{Auto-Correlation}), hay PatchTST~\cite{nie2023time} phân đoạn chuỗi thành các token dạng bản vá đa biến (\textit{subseries patching}).

Tuy nhiên, như đã được chứng minh định lượng trong bài báo của Roy và cộng sự~\cite{roy2025hourly}, các mô hình Transformer phức tạp này thường có xu hướng \textit{underperform} so với các mô hình GBDT phân đoạn theo giờ trên dữ liệu STLF thực tế. Có ba nguyên nhân học thuật cốt lõi giải thích cho hiện tượng này:
\begin{enumerate}
    \item \textbf{Thiếu thiên kiến quy nạp thời gian (\textit{Weak Temporal Inductive Bias}):} Cơ chế tự chú ý xem xét mối liên hệ giữa các điểm thời gian một cách bình đẳng và chỉ dựa vào mã hóa vị trí (\textit{positional encoding}), khiến mô hình dễ bị phân tán bởi các nhiễu tần số cao thay vì ưu tiên các chu kỳ nhật triều và tuần hoàn cố định như mô hình cây phân đoạn.
    \item \textbf{Hiện tượng quá khớp trên dữ liệu ngoại sinh phức tạp:} Khác với ngôn ngữ tự nhiên hay thị giác máy tính, dữ liệu khí tượng chuỗi thời gian chứa tỷ lệ tín hiệu trên nhiễu (\textit{signal-to-noise ratio}) tương đối thấp. Việc tối ưu hóa hàng triệu tham số của Transformer trên chuỗi thời gian 2--3 năm rất dễ dẫn tới hiện tượng ghi nhớ mẫu quá khớp (\textit{memorization}) thay vì học được quy luật tổng quát hóa.
    \item \textbf{Độ trễ và chi phí tính toán:} Chi phí huấn luyện và suy luận của Transformer lớn hơn hàng chục lần so với XGBoost, tạo rào cản lớn khi triển khai thực tế trên các hệ thống giám sát điều độ thời gian thực của các công ty điện lực.
\end{enumerate}

\subsection{Học máy Tích hợp Tri thức Vật lý và Kỹ thuật Stacking Ensemble}
Kỹ thuật tổng quát hóa xếp chồng (\textit{Stacked Generalization} hay \textit{Stacking}) được đề xuất lần đầu bởi Wolpert (1992)~\cite{wolpert1992stacked}, cho phép kết hợp sức mạnh của nhiều mô hình cơ sở thông qua một mô hình tổng hợp cấp cao (\textit{meta-learner}). Trong bài toán STLF, đa số các nghiên cứu Stacking truyền thống chỉ đưa vector dự báo của các mô hình cơ sở $[\hat{y}_1, \dots, \hat{y}_M]$ vào meta-learner. Cách tiếp cận tĩnh này bỏ qua hoàn toàn bối cảnh môi trường vận hành thời gian thực, khiến meta-learner không thể học được quy tắc phân quyền: ví dụ, khi nhiệt độ cực kỳ cao, mô hình nào đáng tin cậy hơn; khi vào ban đêm yên tĩnh, mô hình nào đưa ra dự báo ổn định hơn.

Bên cạnh đó, trong vật lý nhiệt công trình, tiêu chuẩn ASHRAE~\cite{ashrae2021fundamentals} đã chứng minh khái niệm độ-ngày (\textit{Degree Days} bao gồm HDD và CDD) là công cụ toán học tối ưu để xấp xỉ tuyến tính từng đoạn mối quan hệ giữa nhiệt độ môi trường và hành vi tiêu thụ điện của các hệ thống HVAC. Tuy nhiên, việc kết hợp các đặc trưng vật lý này vào mô hình phân đoạn 24 giờ kết hợp kiến trúc Stacking nhận biết ngữ cảnh vẫn là một khoảng trống nghiên cứu chưa từng được khám phá.

\subsection{Các Câu hỏi Nghiên cứu (Research Questions)}
Từ các khoảng trống học thuật trên, đề tài tập trung trả lời ba câu hỏi nghiên cứu cốt lõi:
\begin{itemize}
    \item \textbf{RQ1:} Việc bổ sung các đặc trưng nhiệt động lực học HDD/CDD (ngưỡng $18^\circ\text{C}$) và sóng điều hòa Fourier liên tục vào mô hình 24-XGBoost có cải thiện được độ chính xác và khả năng tổng quát hóa trên tập dữ liệu kiểm tra năm thứ ba so với baseline của Roy et al. (2025) hay không?
    \item \textbf{RQ2:} Mức độ độc lập và tương quan phần dư giữa mô hình tuyến tính ElasticNet (L1/L2 regularization) và mô hình phi tuyến dạng cây XGBoost là bao nhiêu, và mỗi mô hình thể hiện ưu thế vượt trội trong những chế độ vận hành cụ thể nào của lưới điện?
    \item \textbf{RQ3:} Kiến trúc Context-Aware Stacking (CAS) đề xuất, thông qua việc điều hòa meta-learner bằng vector ngữ cảnh khí tượng và thời gian thực, có khả năng khai thác tối đa tính bù trừ của các mô hình cơ sở để tạo nên bước nhảy vọt về hiệu năng so với các mô hình đơn lẻ hay không?
\end{itemize}
